
\documentclass{beamer}

%
% Choose how your presentation looks.
%
% For more themes, color themes and font themes, see:
% http://deic.uab.es/~iblanes/beamer_gallery/index_by_theme.html
%

\mode<presentation>
{
  \usetheme{default}
  \usecolortheme{default}
  \usefonttheme{default}
  \setbeamertemplate{navigation symbols}{}
  \setbeamertemplate{caption}[numbered]
}

\usepackage[english]{babel}
\usepackage[utf8]{inputenc}
\usepackage{amsmath}
\usepackage{amssymb}
\usepackage{booktabs}

\title[Statistics]{Introduction to Statistics}
\author{Guillermo Cabrera}
\institute{Statistics and Data Analysis}
\date{\today}

\begin{document}

% ============================================================
% TITLE
% ============================================================

\begin{frame}
  \titlepage
\end{frame}


% ============================================================
% OUTLINE
% ============================================================

% Uncomment these lines for an automatically generated outline.
%\begin{frame}{Outline}
%  \tableofcontents
%\end{frame}


% ============================================================
% SECTION 1
% ============================================================

\section{Introduction}

\begin{frame}{What is Statistics?}

\begin{itemize}

  \item Statistics is the discipline concerned with the collection, organization, analysis, interpretation and presentation of data.

  \vskip 0.3cm

  \item Its objective is to extract useful information from data and support decision-making under uncertainty.

  \vskip 0.3cm

  \item Statistics can be divided into two major areas:

  \begin{itemize}
    \item Descriptive statistics
    \item Inferential statistics
  \end{itemize}

  \vskip 0.3cm

  \item Statistical analysis provides a framework for understanding patterns, variability and relationships in data.

\end{itemize}

\end{frame}


\begin{frame}{Descriptive and Inferential Statistics}

\begin{itemize}

  \item \textbf{Descriptive statistics} summarizes and describes the information contained in a dataset.

  \vskip 0.4cm

  Examples include:

  \begin{itemize}
    \item Mean
    \item Median
    \item Standard deviation
    \item Percentiles
    \item Tables and graphs
  \end{itemize}

  \vskip 0.4cm

  \item \textbf{Inferential statistics} uses information from a sample to make conclusions about a larger population.

  \vskip 0.4cm

  \item Inferential methods include estimation, hypothesis testing and regression analysis.

\end{itemize}

\end{frame}


% ============================================================
% SECTION 2
% ============================================================

\section{Population and Sample}

\begin{frame}{Population and Sample}

\begin{itemize}

  \item A \textbf{population} is the complete set of individuals, observations or units of interest.

  \vskip 0.4cm

  \item A \textbf{sample} is a subset of the population selected for analysis.

  \vskip 0.4cm

  \item Studying an entire population can be expensive, slow or impossible.

  \vskip 0.4cm

  \item Statistical inference allows us to use information from a sample to learn about the population.

\end{itemize}

\vskip 0.6cm

\begin{center}

Population $\rightarrow$ Sampling $\rightarrow$ Sample $\rightarrow$ Statistical inference

\end{center}

\end{frame}


\begin{frame}{Parameter and Statistic}

\begin{itemize}

  \item A \textbf{parameter} is a numerical characteristic of a population.

  \vskip 0.3cm

  \item A \textbf{statistic} is a numerical characteristic calculated from a sample.

  \vskip 0.4cm

  Examples:

  \begin{itemize}
    \item Population mean: $\mu$
    \item Sample mean: $\bar{x}$
    \item Population variance: $\sigma^2$
    \item Sample variance: $s^2$
  \end{itemize}

  \vskip 0.4cm

  \item The statistic is often used to estimate the unknown population parameter.

\end{itemize}

\end{frame}


% ============================================================
% SECTION 3
% ============================================================

\section{Variables and Data}

\begin{frame}{Statistical Variables}

\begin{itemize}

  \item A \textbf{variable} is a characteristic that can take different values across observations.

  \vskip 0.4cm

  \item Variables can be classified as:

  \begin{itemize}

    \item \textbf{Qualitative variables:} describe categories or attributes.

    \item \textbf{Quantitative variables:} represent numerical quantities.

  \end{itemize}

  \vskip 0.4cm

  \item Quantitative variables can be:

  \begin{itemize}
    \item Discrete
    \item Continuous
  \end{itemize}

\end{itemize}

\end{frame}


\begin{frame}{Types of Variables}

\begin{itemize}

  \item \textbf{Nominal:} categories without an intrinsic order.

  \vskip 0.2cm

  Example: country, gender, industry.

  \vskip 0.4cm

  \item \textbf{Ordinal:} categories with an ordered structure.

  \vskip 0.2cm

  Example: low, medium and high.

  \vskip 0.4cm

  \item \textbf{Discrete:} numerical variables based on counts.

  \vskip 0.2cm

  Example: number of firms.

  \vskip 0.4cm

  \item \textbf{Continuous:} numerical variables that can take values within an interval.

  \vskip 0.2cm

  Example: income, height or temperature.

\end{itemize}

\end{frame}


% ============================================================
% SECTION 4
% ============================================================

\section{Measures of Central Tendency}

\begin{frame}{Measures of Central Tendency}

\begin{itemize}

  \item Measures of central tendency describe the central or typical value of a dataset.

  \vskip 0.4cm

  \item The three most common measures are:

  \begin{itemize}
    \item Mean
    \item Median
    \item Mode
  \end{itemize}

  \vskip 0.4cm

  \item The appropriate measure depends on the distribution and characteristics of the data.

\end{itemize}

\end{frame}


\begin{frame}{Arithmetic Mean}

The arithmetic mean is calculated as:

\vskip 0.5cm

\[
\bar{x} =
\frac{1}{n}
\sum_{i=1}^{n} x_i
\]

\vskip 0.5cm

\begin{itemize}

  \item $x_i$ represents each observation.

  \item $n$ is the number of observations.

  \item The mean uses every observation in the dataset.

  \item It is sensitive to extreme values.

\end{itemize}

\vskip 0.3cm

For example:

\[
2,\;4,\;6,\;8,\;10
\]

\[
\bar{x} = \frac{2+4+6+8+10}{5}=6
\]

\end{frame}


\begin{frame}{Median and Mode}

\begin{itemize}

  \item The \textbf{median} is the value that separates the ordered data into two equal parts.

  \vskip 0.3cm

  \item The \textbf{mode} is the most frequently occurring value.

\end{itemize}

\vskip 0.5cm

For the dataset:

\[
2,\;3,\;3,\;4,\;5,\;8,\;10
\]

\vskip 0.3cm

\[
\text{Mean} = 5
\]

\[
\text{Median} = 4
\]

\[
\text{Mode} = 3
\]

\vskip 0.4cm

The median is generally more robust than the mean when extreme observations are present.

\end{frame}


% ============================================================
% SECTION 5
% ============================================================

\section{Measures of Dispersion}

\begin{frame}{Measures of Dispersion}

\begin{itemize}

  \item Measures of dispersion describe how spread out the observations are.

  \vskip 0.4cm

  Important measures include:

  \begin{itemize}
    \item Range
    \item Variance
    \item Standard deviation
    \item Interquartile range
  \end{itemize}

  \vskip 0.4cm

  \item Two datasets can have the same mean but very different levels of variability.

\end{itemize}

\end{frame}


\begin{frame}{Variance}

The sample variance is defined as:

\vskip 0.5cm

\[
s^2 =
\frac{1}{n-1}
\sum_{i=1}^{n}
(x_i-\bar{x})^2
\]

\vskip 0.5cm

\begin{itemize}

  \item It measures the average squared deviation from the sample mean.

  \item Squaring the deviations prevents positive and negative differences from cancelling each other.

  \item The denominator $n-1$ provides the usual unbiased estimator of population variance under standard random-sampling assumptions.

\end{itemize}

\end{frame}


\begin{frame}{Standard Deviation}

The standard deviation is the square root of the variance:

\vskip 0.5cm

\[
s =
\sqrt{
\frac{1}{n-1}
\sum_{i=1}^{n}
(x_i-\bar{x})^2
}
\]

\vskip 0.5cm

\begin{itemize}

  \item Standard deviation is expressed in the same units as the original variable.

  \item A small standard deviation indicates observations are concentrated around the mean.

  \item A large standard deviation indicates greater dispersion.

\end{itemize}

\end{frame}


\begin{frame}{Coefficient of Variation}

The coefficient of variation measures relative dispersion:

\vskip 0.5cm

\[
CV =
\frac{s}{\bar{x}}
\]

\vskip 0.2cm

or, expressed as a percentage,

\[
CV =
\frac{s}{\bar{x}}\times100
\]

\vskip 0.5cm

\begin{itemize}

  \item It is useful for comparing variability across variables with different scales.

  \item It is particularly useful when the mean is positive and meaningfully different from zero.

\end{itemize}

\end{frame}


% ============================================================
% SECTION 6
% ============================================================

\section{Distribution of Data}

\begin{frame}{Frequency Distribution}

\begin{itemize}

  \item A frequency distribution describes how often different values or intervals occur.

  \vskip 0.4cm

  \item Frequencies can be represented using:

  \begin{itemize}
    \item Frequency tables
    \item Histograms
    \item Bar charts
    \item Density plots
  \end{itemize}

  \vskip 0.4cm

  \item Visualization helps identify concentration, dispersion, skewness and unusual observations.

\end{itemize}

\end{frame}


\begin{frame}{Normal Distribution}

The normal distribution is one of the most important probability distributions in statistics.

\vskip 0.5cm

Its probability density function is:

\[
f(x)=
\frac{1}{\sigma\sqrt{2\pi}}
e^{-\frac{1}{2}
\left(
\frac{x-\mu}{\sigma}
\right)^2}
\]

\vskip 0.5cm

\begin{itemize}

  \item It is symmetric around the mean $\mu$.

  \item The mean, median and mode are equal.

  \item The parameter $\sigma$ determines the dispersion of the distribution.

\end{itemize}

\end{frame}


\begin{frame}{The Empirical Rule}

For a normally distributed variable:

\vskip 0.5cm

\begin{itemize}

  \item Approximately $68\%$ of observations lie within $\mu \pm 1\sigma$.

  \item Approximately $95\%$ lie within $\mu \pm 2\sigma$.

  \item Approximately $99.7\%$ lie within $\mu \pm 3\sigma$.

\end{itemize}

\vskip 0.7cm

\[
68\% \quad \rightarrow \quad 1\sigma
\]

\[
95\% \quad \rightarrow \quad 2\sigma
\]

\[
99.7\% \quad \rightarrow \quad 3\sigma
\]

\end{frame}


% ============================================================
% SECTION 7
% ============================================================

\section{Probability}

\begin{frame}{Probability}

\begin{itemize}

  \item Probability quantifies uncertainty about the occurrence of an event.

  \vskip 0.4cm

  \item The probability of an event $A$ satisfies:

\[
0 \leq P(A) \leq 1
\]

  \vskip 0.4cm

  \item A probability of $0$ means the event is impossible.

  \item A probability of $1$ means the event is certain.

\end{itemize}

\end{frame}


\begin{frame}{Conditional Probability}

Conditional probability measures the probability of an event given that another event has occurred.

\vskip 0.5cm

\[
P(A|B)
=
\frac{P(A\cap B)}
{P(B)}
\]

\vskip 0.5cm

\begin{itemize}

  \item $P(A|B)$ is the probability of $A$ given $B$.

  \item $P(A\cap B)$ is the probability that both events occur.

  \item $P(B)$ is the probability of event $B$.

\end{itemize}

\end{frame}


% ============================================================
% SECTION 8
% ============================================================

\section{Correlation}

\begin{frame}{Correlation}

\begin{itemize}

  \item Correlation measures the strength and direction of a linear relationship between two quantitative variables.

  \vskip 0.5cm

  \item The Pearson correlation coefficient is:

\[
r =
\frac{
\sum_{i=1}^{n}
(x_i-\bar{x})(y_i-\bar{y})
}{
\sqrt{
\sum_{i=1}^{n}(x_i-\bar{x})^2
\sum_{i=1}^{n}(y_i-\bar{y})^2
}
}
\]

  \vskip 0.4cm

  \item The coefficient satisfies:

\[
-1 \leq r \leq 1
\]

\end{itemize}

\end{frame}


\begin{frame}{Interpreting Correlation}

\begin{itemize}

  \item $r \approx 1$: strong positive linear association.

  \vskip 0.3cm

  \item $r \approx -1$: strong negative linear association.

  \vskip 0.3cm

  \item $r \approx 0$: weak or no linear association.

  \vskip 0.5cm

  \item Correlation does not by itself establish causality.

  \vskip 0.3cm

  \item A high correlation can arise from confounding variables, common trends or other mechanisms.

\end{itemize}

\end{frame}


% ============================================================
% SECTION 9
% ============================================================

\section{Sampling}

\begin{frame}{Sampling Methods}

\begin{itemize}

  \item Sampling determines which observations are selected from a population.

  \vskip 0.4cm

  Common probability sampling methods include:

  \begin{itemize}
    \item Simple random sampling
    \item Systematic sampling
    \item Stratified sampling
    \item Cluster sampling
  \end{itemize}

  \vskip 0.4cm

  \item The sampling design affects the representativeness and precision of statistical estimates.

\end{itemize}

\end{frame}


\begin{frame}{Sampling Error}

\begin{itemize}

  \item A sample generally does not reproduce the population perfectly.

  \vskip 0.4cm

  \item \textbf{Sampling error} is the random difference between a sample estimate and the corresponding population parameter.

  \vskip 0.4cm

  \item Increasing the sample size generally reduces sampling variability.

  \vskip 0.4cm

  \item However, increasing sample size does not automatically eliminate systematic bias.

\end{itemize}

\end{frame}


% ============================================================
% SECTION 10
% ============================================================

\section{Statistical Inference}

\begin{frame}{Confidence Intervals}

A confidence interval provides a range of plausible values for a population parameter.

\vskip 0.5cm

For a population mean under standard conditions:

\[
\bar{x}
\pm
z_{\alpha/2}
\frac{\sigma}{\sqrt{n}}
\]

\vskip 0.5cm

\begin{itemize}

  \item $\bar{x}$ is the sample mean.

  \item $\sigma$ is the population standard deviation.

  \item $n$ is the sample size.

  \item $z_{\alpha/2}$ is the appropriate critical value.

\end{itemize}

\end{frame}


\begin{frame}{Hypothesis Testing}

\begin{itemize}

  \item Hypothesis testing provides a formal framework for evaluating claims about a population.

  \vskip 0.4cm

  \item The procedure usually involves:

  \begin{enumerate}
    \item Formulating the null hypothesis $H_0$.
    \item Formulating the alternative hypothesis $H_1$.
    \item Choosing a significance level $\alpha$.
    \item Calculating a test statistic.
    \item Obtaining a $p$-value or comparing with a critical value.
    \item Making a statistical decision.
  \end{enumerate}

\end{itemize}

\end{frame}


\begin{frame}{Statistical Significance}

\begin{itemize}

  \item The $p$-value measures how incompatible the observed data are with the null hypothesis, under the assumptions of the statistical test.

  \vskip 0.4cm

  \item A common significance level is:

\[
\alpha = 0.05
\]

  \vskip 0.4cm

  \item If $p < \alpha$, the null hypothesis is rejected under the chosen testing rule.

  \vskip 0.4cm

  \item Statistical significance does not necessarily imply economic, practical or causal significance.

\end{itemize}

\end{frame}


% ============================================================
% SECTION 11
% ============================================================

\section{Regression}

\begin{frame}{Simple Linear Regression}

Simple linear regression models the conditional mean of a dependent variable as a function of one explanatory variable:

\vskip 0.5cm

\[
Y_i =
\beta_0 +
\beta_1 X_i +
\varepsilon_i
\]

\vskip 0.5cm

\begin{itemize}

  \item $Y_i$: dependent variable.

  \item $X_i$: explanatory variable.

  \item $\beta_0$: intercept.

  \item $\beta_1$: slope coefficient.

  \item $\varepsilon_i$: error term.

\end{itemize}

\end{frame}


\begin{frame}{Interpreting the Regression Coefficient}

For the model:

\[
Y_i =
\beta_0 +
\beta_1X_i+
\varepsilon_i
\]

\vskip 0.5cm

\begin{itemize}

  \item $\beta_0$ represents the expected value of $Y$ when $X=0$.

  \vskip 0.3cm

  \item $\beta_1$ represents the expected change in $Y$ associated with a one-unit increase in $X$, conditional on the model specification.

  \vskip 0.4cm

  \item Interpretation depends on the measurement units, functional form and assumptions of the model.

\end{itemize}

\end{frame}


\begin{frame}{Coefficient of Determination}

The coefficient of determination is:

\vskip 0.5cm

\[
R^2 =
1 -
\frac{SS_{Residual}}
{SS_{Total}}
\]

\vskip 0.5cm

\begin{itemize}

  \item $R^2$ measures the proportion of variation in the dependent variable explained by the fitted model relative to a constant-only benchmark.

  \vskip 0.3cm

  \item Higher $R^2$ does not necessarily imply a better causal model.

  \vskip 0.3cm

  \item Model quality should also consider assumptions, specification, prediction error and out-of-sample performance.

\end{itemize}

\end{frame}


% ============================================================
% SECTION 12
% ============================================================

\section{Conclusion}

\begin{frame}{Key Statistical Concepts}

\begin{itemize}

  \item Statistics provides tools for understanding data and uncertainty.

  \vskip 0.3cm

  \item Descriptive statistics summarize observed data.

  \vskip 0.3cm

  \item Probability provides the mathematical framework for uncertainty.

  \vskip 0.3cm

  \item Sampling allows us to study populations using subsets of observations.

  \vskip 0.3cm

  \item Inferential statistics allows estimation and hypothesis testing.

  \vskip 0.3cm

  \item Correlation measures association, while regression provides a framework for modeling conditional relationships.

  \vskip 0.3cm

  \item Statistical results must always be interpreted considering assumptions, uncertainty and potential sources of bias.

\end{itemize}

\end{frame}


\begin{frame}{Final Perspective}

\begin{center}

\Large
``Statistics is not only about calculating numbers.''

\vskip 0.8cm

\normalsize

It is about understanding variation,\\
quantifying uncertainty,\\
and making informed decisions from data.

\end{center}

\end{frame}


\end{document}
```
